\par\clearpage
\normalmarginpar
\scanpage{3}
\noindent\hfill\margin{\unsure{zurück} \unsure{an|au} \unsure{·}\tikz[baseline=(g.base)]{\node[inner sep=0pt] (g) {G};\draw ([xshift=-1pt,yshift=-4pt]g.south west)--([xshift=4pt,yshift=1pt]g.south east);}}\par

\noindent SGA 7

\noindent D\quad $\underline{\overline{\mathrm{X}}}$\quad \unsure{Intersections|Intersection} \tikz[baseline]{\draw[line cap=round] (0,.035)--(.065,.04) (.12,.04)--(.151,.043) (.211,.034)--(.295,.057);}

\noindent D\quad $\underline{\overline{\mathrm{XI}}}$\quad Int. \unsure{complètes|completes}

\noindent D\quad $\underline{\overline{\mathrm{XII}}}$\quad \unsure{Quadriques|Quadrique}

\noindent D\quad $\underline{\overline{\mathrm{XIII}}}$\quad Points \unsure{quadr.} \unsure{ordinaires}. Pic-\unsure{Lef|lef} faible

\noindent D\quad $\underline{\overline{\mathrm{XIV}}}$\quad Formule \unsure{Milnor}

\diagram{\(\begin{array}{@{}ll@{}}
\text{D}&\underline{\overline{\mathrm{XV}}}\\
\text{D}&\underline{\overline{\mathrm{XVI}}}
\end{array}
\left\{\begin{array}{l}
\text{Picard\unsure{'}-Lefschetz}\\
\text{Picard-Lefschetz}
\end{array}\right.
\;\tikz[baseline]{\draw[gray,thin] (0,.86) .. controls (.035,.62) and (-.045,.06) .. (.013,-.23);}
\)}

\noindent K\quad
$\begin{array}{@{}c@{}}
\underline{\overline{\mathrm{XVII}}}\\[-2pt]
\underline{\overline{\mathrm{XVIII}}}
\end{array}\,\Bigg|$\quad
\unsure{Pinceaux} de Lefschetz~: \uline{existence} \tikz[baseline]{\draw[gray,thin] (.18,.62) .. controls (.19,.36) and (.08,.18) .. (.07,-.09) -- (.11,-.17);\draw[gray,thin] (.20,.58) .. controls (.15,.25) and (.01,-.02) .. (-.01,-.20) -- (-.035,-.16);\draw[line cap=round] (-.06,0)--(.17,.34) (-.04,-.015)--(.18,.31) (.03,.013)--(.25,.35);} \uline{\unsure{étude|études} \unsure{cohomologique|cohomologiques}}

\diagram{D\quad $\underline{\overline{\deleted{\unsure{\mathrm{XVIII}}}\mathrm{XIX}}}$\quad
\uuline{Th de Noether}
$\left\{\begin{array}{l}
\text{via \unsure{Hodge}-Deligne}\\
\text{via \uuline{\underline{\text{Katz}}}}
\end{array}\right.$}

\diagram{\begin{tabular}{@{}c@{\ }p{.84\linewidth}@{}}
$\mathrm{K}\left\{\rule{0pt}{12ex}\right.$ &
\begin{minipage}[c]{\linewidth}
\noindent $\deleted{\underline{\overline{\unsure{\mathrm{XIX}|\mathrm{XX}}}}}\;
\underline{\overline{\mathrm{XX}}}\deleted{\unsure{\mathrm{I}}}$\quad \unsure{Th.|Thr.} de Griffiths.

\noindent $\underline{\overline{\mathrm{XXI}}}$\quad \deleted{\unsure{Théorème|Théor}} Niveau de la coh. des \unsure{intersections} \unsure{complètes|completes}

\noindent $\underline{\overline{\mathrm{XXII}}}$\quad \deleted{\unsure{Form.|Fon.}} Calcul de la \unsure{fonction} $\deleted{\unsure{z}}\zeta$ mod $p$.
\end{minipage}
\end{tabular}}

\figure{\begin{tikzpicture}\draw (0,0) .. controls (.29\linewidth,.45) and (.76\linewidth,.57) .. (.98\linewidth,.46);\end{tikzpicture}}

\noindent Exposés \unsure{transcendants} (Lê Dũng Tráng ?)

\[
\underline{\mathbb{F}_{q}\qquad \#\qquad \mathrm{card}}\quad\unsure{.}
\]

% archive sheet number: 2
